%% FullCourse_10Lecs -- Lecture 9, Topic 9.1b: The Agentic-AI Landscape (mid-2026)
%% New deck broadening the lecture from the terminal-agent spine to the
%% mid-2026 agentic-AI landscape. Facts web-verified for June 2, 2026.
%% Tagging convention: [solid] = primary/citable; [hype] = vendor self-report.

%% Shared preamble for the FullCourse_10Lecs topic decks.
%%
%% Each topic .tex \input's this file, then defines its own \title, \date
%% override (if any), \graphicspath, and \begin{document}.
%%
%% Reconstructed verbatim from the inlined copy preserved in
%% Lec10_Case_Studies/Lec10_T8_Case6_DeepRamsey.tex (lines 23-190), which is
%% explicitly annotated there as "from ../shared_preamble.tex, copied verbatim".
%% Base preamble matches MiniCourse_8hr/Slides/shared_preamble.tex; this version
%% additionally provides \sectiondivider, the conceptbox/cautionbox/examplebox/
%% takeawaybox tcolorboxes, and \compacttables used across the split decks.

\documentclass[10pt,english,aspectratio=169]{beamer}

\usepackage{ae,aecompl}
\usepackage[T1]{fontenc}
\usepackage{textcomp}
\usepackage[utf8]{inputenc}
\usepackage{babel}
\usepackage{datetime2}

\setcounter{secnumdepth}{3}
\setcounter{tocdepth}{3}

\usepackage{amsmath, amssymb, amsthm, graphicx}
\usepackage{booktabs, multirow}
\usepackage{tikz}
\usepackage{pgfplots}
\pgfplotsset{compat=1.16}
\usepackage[most]{tcolorbox}
\tcbuselibrary{skins, breakable, theorems}
\usepackage{listings}
\usepackage{xcolor}

\lstset{
  basicstyle=\ttfamily\small,
  keywordstyle=\color{blue!80!black}\bfseries,
  commentstyle=\color{green!50!black}\itshape,
  stringstyle=\color{orange!80!black},
  backgroundcolor=\color{gray!8},
  frame=single,
  rulecolor=\color{gray!40},
  breaklines=true,
  showstringspaces=false,
  numbers=none,
}

\lstdefinestyle{pythonstyle}{
    language=Python,
    basicstyle=\ttfamily\footnotesize,
    keywordstyle=\color{blue!80!black}\bfseries,
    commentstyle=\color{green!50!black}\itshape,
    stringstyle=\color{orange!80!black},
    frame=single,
    breaklines=true,
    showstringspaces=false,
    numbers=left,
    numbersep=5pt,
    numberstyle=\tiny\color{gray},
    backgroundcolor=\color{gray!8},
    rulecolor=\color{gray!40}
}

\ifx\hypersetup\undefined
  \AtBeginDocument{%
    \hypersetup{bookmarks=false,breaklinks=true,pdfborder={0 0 0},colorlinks=true,
      linkcolor=DarkText,urlcolor=RedTitle}%
  }
\else
  \hypersetup{bookmarks=false,breaklinks=true,pdfborder={0 0 0},colorlinks=true,
    linkcolor=DarkText,urlcolor=RedTitle}%
\fi

\makeatletter
\newcommand\makebeamertitle{%
  \begin{frame}[plain]
    \vspace*{1.0cm}
    \centering
    {\usebeamerfont{title}\usebeamercolor[fg]{title}\inserttitle\par}
    \vspace{1.0cm}
    {\usebeamerfont{author}\usebeamercolor[fg]{author}\insertauthor\par}
    \vspace{0.2cm}
    {\usebeamerfont{date}\usebeamercolor[fg]{date}\insertdate\par}
  \end{frame}}%
\AtBeginDocument{%
  \let\origtableofcontents=\tableofcontents
  \def\tableofcontents{\@ifnextchar[{\origtableofcontents}{\gobbletableofcontents}}
  \def\gobbletableofcontents#1{\origtableofcontents}%
}

\usetheme{metropolis}
\usefonttheme{professionalfonts}

\definecolor{LightGray}{RGB}{230,230,230}
\definecolor{RedTitle}{RGB}{0,0,200}
\definecolor{VeryLightGray}{RGB}{245,245,245}
\definecolor{DarkText}{RGB}{0,0,0}

\setbeamercolor{background canvas}{bg=white}
\setbeamercolor{title}{fg=RedTitle}
\setbeamercolor{author}{fg=DarkText}
\setbeamercolor{date}{fg=DarkText}
\setbeamercolor{frametitle}{bg=VeryLightGray,fg=DarkText}
\setbeamercolor{normal text}{fg=DarkText}
\setbeamerfont{title}{size=\large}

\setbeamersize{text margin left=5mm, text margin right=5mm}
\addtobeamertemplate{frametitle}{}{\vspace{-0.5em}}
\newcommand{\reducevspace}{\vspace{-0.5cm}}

% Shared section divider and box styles for split topic decks.
\newcommand{\sectiondivider}[2]{%
\begin{frame}[plain,noframenumbering]
\begin{center}
\vfill
{\Large\textbf{#1}\par}
\vspace{0.35cm}
{\normalsize #2\par}
\vfill
\end{center}
\end{frame}
}

\newtcolorbox{conceptbox}[1][]{
  colback=blue!3,
  colframe=RedTitle!55,
  boxrule=0.35mm,
  arc=1mm,
  left=2mm,
  right=2mm,
  top=1mm,
  bottom=1mm,
  #1
}
\newtcolorbox{cautionbox}[1][]{
  colback=red!3,
  colframe=red!50!black,
  boxrule=0.35mm,
  arc=1mm,
  left=2mm,
  right=2mm,
  top=1mm,
  bottom=1mm,
  #1
}
\newtcolorbox{examplebox}[1][]{
  colback=green!3,
  colframe=green!45!black,
  boxrule=0.35mm,
  arc=1mm,
  left=2mm,
  right=2mm,
  top=1mm,
  bottom=1mm,
  #1
}
\newtcolorbox{takeawaybox}[1][]{
  colback=VeryLightGray,
  colframe=RedTitle!55,
  boxrule=0.35mm,
  arc=1mm,
  left=2mm,
  right=2mm,
  top=1mm,
  bottom=1mm,
  #1
}
\newcommand{\compacttables}{\renewcommand{\arraystretch}{1.15}}

\setbeamertemplate{navigation symbols}{}
\setbeamertemplate{footline}{%
  \hfill%
  \usebeamercolor[fg]{page number in head/foot}%
  \usebeamerfont{page number in head/foot}%
  \insertframenumber\,/\,\inserttotalframenumber\kern1em\vskip2pt%
}

\makeatother

\author{Zhigang Feng}
% "date with time": compile timestamp so each build is identifiable.
\date{\today\ \DTMcurrenttime}


% Suppress redundant auto-generated metropolis section page; \sectiondivider frames are the dividers we keep.
\metroset{sectionpage=none}

\graphicspath{{./pic/}{./figures/}{../../MiniCourse_8hr/Slides/Module4_Agentic_AI/pic/}{../}}


\usetikzlibrary{positioning,arrows.meta,shapes,calc}

\title[AI for Econ Research]{AI for Economic Research: Dynamic Models, Language, and Agents\\
Lecture 9.1b: The Agentic-AI Landscape (mid-2026)}

\begin{document}

\makebeamertitle

%=============================================================================
\begin{frame}{Topic 9.1b Agenda}
\textbf{T1 framed the terminal-agent workflow. This interlude widens the lens to the broader agentic-AI landscape an economist now operates inside.}

\vspace{0.2cm}

\begin{enumerate}
  \item \textbf{The interoperability stack} --- MCP, A2A, AGNTCY/ACP, neutral governance
  \item \textbf{Orchestration patterns} --- frameworks, context engineering, the multi-agent cost-benefit
  \item \textbf{Reasoning models \& inference-time compute} --- think budgets, measured time horizons
  \item \textbf{Autonomous coding \& computer-use agents} --- scheduled runtimes, sandboxing
  \item \textbf{Benchmarks \& measurement validity} --- ties to the course's measurement theme
  \item \textbf{Safety \& reproducibility} --- untrusted inputs, chain-of-custody
\end{enumerate}

\vspace{0.15cm}
\footnotesize\textit{Reading convention: \textbf{[solid]} $=$ primary or citable; \textbf{[hype]} $=$ vendor self-report (counter-source given).}

\vspace{0.05cm}
{\scriptsize \textit{Facts current as of June 2026. Model versions are stated with release dates; the field moves monthly --- recheck before teaching.}}
\end{frame}

%=============================================================================
\section{The Interoperability Stack}
%===========================================================

\sectiondivider{The Interoperability Stack}{Compete on models, cooperate on plumbing}

\begin{frame}{Three Layers of Agent Interoperability}
\textbf{The 2025--26 story is standard-setting: vendors compete on models but increasingly cooperate on the wiring between agents and tools.}

\vspace{0.18cm}

\begin{center}
\footnotesize
\renewcommand{\arraystretch}{1.2}
\begin{tabular}{|p{2.6cm}|p{3.0cm}|p{4.9cm}|}
\hline
\textbf{Layer} & \textbf{Protocol} & \textbf{What it standardizes} \\ \hline
Agent $\leftrightarrow$ tool & \textbf{MCP} (Model Context Protocol) & how an agent discovers and calls external tools/data sources \\ \hline
Agent $\leftrightarrow$ agent & \textbf{A2A} (Agent-to-Agent) & how independent agents delegate tasks to one another \\ \hline
Discovery / identity & \textbf{AGNTCY / ACP} & agent directories, capability cards, identity and trust \\ \hline
\end{tabular}
\end{center}

\vspace{0.18cm}

\begin{itemize}\footnotesize
    \item MCP was introduced by Anthropic (Nov 2024); it is now \emph{multi-vendor}, not proprietary.
    \item \textbf{OpenAI adopted MCP} at DevDay (Oct 6, 2025); Google and Microsoft tooling support it too. \textbf{[solid]}
    \item Governance moved to a neutral home: the \textbf{Linux Foundation Agentic AI Foundation (AAIF)}, formed Dec 9, 2025, now hosts MCP and related projects. \textbf{[solid]}
\end{itemize}
\end{frame}

\begin{frame}{Why Neutral Governance Is a Network-Effects Story}
\textbf{For economists, the interoperability stack is a textbook standard-setting / network-effects case.}

\vspace{0.22cm}

\begin{itemize}
    \item A shared protocol lowers switching costs: a tool wrapped once in MCP works across Claude, ChatGPT, Gemini, and open agents
    \item Network effects make the standard more valuable as adoption grows --- which is why a single-vendor protocol faces adoption resistance
    \item Moving MCP under the \textbf{Linux Foundation AAIF} (Dec 9, 2025) is the classic move: hand the ``plumbing'' to a neutral body so rivals will build on it
    \item Competition then concentrates where rents remain: model quality, latency, and price --- not the connectors
\end{itemize}

\vspace{0.22cm}

\begin{takeawaybox}
\footnotesize\textbf{Reproducibility note.} Protocols version. Pin the \emph{spec version} you used (e.g.\ MCP spec date, A2A version) alongside the model version and date --- a tool call that worked under one spec revision may behave differently under the next.
\end{takeawaybox}
\end{frame}

%=============================================================================
\section{Orchestration Patterns}
%===========================================================

\sectiondivider{Orchestration Patterns}{How multi-step agent systems are actually built}

\begin{frame}{Four Orchestration Patterns}
\textbf{Production agent systems are built from a small set of recurring control structures.}

\vspace{0.18cm}

\begin{center}
\footnotesize
\renewcommand{\arraystretch}{1.2}
\begin{tabular}{|p{3.0cm}|p{7.4cm}|}
\hline
\textbf{Pattern} & \textbf{Shape} \\ \hline
Orchestrator--worker & a lead agent decomposes a task and dispatches workers; \emph{dominant in production} \\ \hline
Planner--executor & one component plans, another executes step by step \\ \hline
Reflection & the agent critiques and revises its own output before returning \\ \hline
Handoff & control passes between specialized agents (e.g.\ triage $\to$ specialist) \\ \hline
\end{tabular}
\end{center}

\vspace{0.16cm}

\footnotesize\textbf{The terminal-agent workflow you learned in T1--T5 is an orchestrator--worker system}: the main agent orchestrates; \texttt{code-reviewer}, \texttt{domain-reviewer}, and \texttt{verifier} sub-agents are the workers. The landscape just adds vocabulary and tooling around the same loop.
\end{frame}

\begin{frame}{Frameworks Reached 1.0 (2025--26)}
\textbf{The orchestration layer matured from research demos to versioned, supported frameworks.}

\vspace{0.16cm}

\begin{center}
\footnotesize
\renewcommand{\arraystretch}{1.15}
\begin{tabular}{|p{4.6cm}|p{3.0cm}|p{2.6cm}|}
\hline
\textbf{Framework} & \textbf{Milestone} & \textbf{Vendor} \\ \hline
LangGraph 1.0 & Oct 2025 & LangChain \\ \hline
OpenAI Agents SDK ($+$ AgentKit) & SDK Mar 2025 & OpenAI \\ \hline
Google Agent Dev. Kit (ADK) & v1.0 May 2025 & Google \\ \hline
Microsoft Agent Framework & 1.0 Apr 2026 & Microsoft \\ \hline
CrewAI & 1.0 & CrewAI \\ \hline
Anthropic \textbf{Claude Agent SDK} & Sep 2025 & Anthropic \\ \hline
\end{tabular}
\end{center}

\vspace{0.12cm}
\footnotesize\textbf{Naming note.} What was the \emph{Claude Code SDK} was renamed the \textbf{Claude Agent SDK} on Sep 29, 2025, reflecting generalization beyond coding. \textbf{[solid]}

\vspace{0.05cm}
\textbf{Durable lesson:} the patterns outlive the framework. Learn orchestrator--worker once; map it onto whichever SDK your institution adopts.
\end{frame}

\begin{frame}{The Multi-Agent Cost--Benefit Decision}
\textbf{More agents is not free. Treat it as marginal value vs.\ marginal cost.}

\vspace{0.2cm}

\begin{columns}[T]
\begin{column}{0.56\textwidth}
\begin{itemize}
    \item Anthropic reports its multi-agent research system beat a single-agent baseline by \textbf{+90.2\%} on its internal research-eval\ldots
    \item \ldots at roughly \textbf{15$\times$ the token cost}
    \item So multi-agent pays off when the task is broad, parallelizable, and the answer's value is high relative to tokens
    \item It is wasteful for narrow, sequential, or cheap tasks
\end{itemize}

\vspace{0.12cm}
{\scriptsize \textbf{[solid, label: Anthropic-internal-eval]} --- a vendor benchmark on a self-defined research task, not an external suite. Read the $+90.2\%$ as directional, not portable.}
\end{column}
\begin{column}{0.42\textwidth}
\begin{takeawaybox}
\footnotesize\textbf{Economist framing.} You are choosing a point where marginal research value $=$ marginal token (and latency, and debugging) cost. The Aiyagari sub-agent review in T4 was justified on \emph{wall-clock}, not because ``more agents are better.''
\end{takeawaybox}
\end{column}
\end{columns}
\end{frame}

\begin{frame}{Context Engineering Displaced Prompt Engineering}
\textbf{The 2025--26 shift: managing what is \emph{in} the context window now matters more than wording a clever prompt.}

\vspace{0.2cm}

\begin{itemize}
    \item \textbf{Compaction.} Summarize and clear stale turns so the working set stays in the high-recall part of the window (the ``lost in the middle'' problem from Lec07 T3b).
    \item \textbf{External memory.} Persist state to files (\texttt{NOTES.md}, \texttt{progress.md}, \texttt{CLAUDE.md}) rather than relying on the chat history.
    \item \textbf{Sub-agents with clean context windows.} Hand a worker a fresh, focused context instead of the orchestrator's accumulated noise.
\end{itemize}

\vspace{0.2cm}

\textbf{Why this matters for research:} it is exactly the discipline T4--T5 already taught --- file-driven workflow, progress logs, focused reviewers. The landscape gave it a name; the habit was already correct.
\end{frame}

%=============================================================================
\section{Reasoning Models \& Inference-Time Compute}
%===========================================================

\sectiondivider{Reasoning Models}{Tool use in the reasoning loop, and how long agents actually run}

\begin{frame}{Reasoning Models and Think-Budget Control}
\textbf{Frontier models now spend variable inference-time compute, and increasingly call tools \emph{inside} the reasoning loop.}

\vspace{0.16cm}

\begin{itemize}\footnotesize
    \item \textbf{Tool use in the loop.} OpenAI o3 / o4-mini (Apr 2025) interleave tool calls with chain-of-thought rather than reasoning first and acting after.
    \item \textbf{Think-budget routers.} GPT-5 (Aug 2025) $\to$ GPT-5.5 (Apr 2026) route between fast and deep modes per query.
    \item \textbf{Toggleable extended thinking + configurable token budget.} \textbf{Claude Opus 4.5} (Nov 2025) lets you turn extended thinking on/off and cap the thinking-token budget; it was the first model past \textbf{80\% on SWE-bench Verified}, and Opus pricing was cut to \$5/\$25 per Mtok (input/output).
    \item \textbf{Parallel deep reasoning.} Gemini ``Deep Think'' explores multiple reasoning paths in parallel.
\end{itemize}

\vspace{0.12cm}
\begin{cautionbox}
\footnotesize\textbf{Benchmark caveat (carry this everywhere).} The 80\% SWE-bench Verified figure is real but the metric is contaminated/saturating: OpenAI \emph{stopped reporting} SWE-bench Verified in Feb 2026. Pair any Verified headline with SWE-bench Pro ($\sim$55--59\% at the frontier, held-out). \textbf{[solid]}
\end{cautionbox}
\end{frame}

\begin{frame}{How Long Can an Agent Actually Work? Measure, Don't Trust Marketing}
\textbf{Vendors say ``works for hours'' or ``runs for days.'' METR gives a \emph{measured} alternative.}

\vspace{0.18cm}

\begin{columns}[T]
\begin{column}{0.56\textwidth}
\begin{itemize}\footnotesize
    \item METR's \emph{task-length-doubling} result: the task length an agent completes at \textbf{50\% reliability} has doubled roughly every \textbf{7 months}, accelerating to $\sim$\textbf{4 months} in 2024--25. \textbf{[solid]}
    \item Time Horizon 1.1 (Jan 2026) pegs \textbf{Opus 4.5 $\approx$ 320 min} and \textbf{GPT-5 $\approx$ 214 min} (50\%-reliability horizon).
    \item Use these measured minutes instead of vendor ``multi-hour/multi-day autonomy.''
\end{itemize}
\end{column}
\begin{column}{0.42\textwidth}
\begin{cautionbox}
\footnotesize\textbf{Metric fragility.} The number is tied to a \textbf{50\% reliability} threshold and a specific task suite; the horizon at 80\% reliability is much shorter, and it shifts across suites. A measured-but-fragile number still beats a vendor claim. \\[2pt]
\textbf{[hype counterpoint]} product pages advertising ``multi-hour'' or ``multi-day'' autonomy --- no reliability threshold stated.
\end{cautionbox}
\end{column}
\end{columns}

\vspace{0.12cm}
\footnotesize\textbf{Research translation:} the right question is not ``can it run for hours?'' but ``at what reliability, on what task distribution, validated how?''
\end{frame}

%=============================================================================
\section{Autonomous Coding \& Computer-Use Agents}
%===========================================================

\sectiondivider{Autonomous Coding \& Computer Use}{From interactive assistants to scheduled, sandboxed runtimes}

\begin{frame}{Coding Agents Became Scheduled, Autonomous Runtimes}
\textbf{In 2025--26, coding agents shifted from interactive assistants to background runtimes with persistent state.}

\vspace{0.16cm}

\begin{itemize}\footnotesize
    \item \textbf{OpenAI Codex ``Goal Mode''} became the default (May 2026), with persistent cross-session state --- the agent resumes a goal rather than starting fresh each chat.
    \item \textbf{Anthropic ``Code with Claude'' (May 6, 2026)} reframed Claude Code as an \emph{orchestrator}:
    \begin{itemize}\scriptsize
        \item \textbf{Auto Mode} with prompt-injection / destructive-action classifiers $+$ human approval gates
        \item \textbf{Routines} triggered by cron or webhook
        \item \textbf{Worktrees} for parallel branches
        \item sandboxed \textbf{Managed Agents} with scoped credentials
    \end{itemize}
\end{itemize}

\vspace{0.12cm}
\begin{takeawaybox}
\footnotesize\textbf{Lesson for research infra.} Before an agent runs unattended on real data or repos, configure sandboxing, credential scoping, and approval gates first. ``Autonomous'' is a permissions decision, not a model capability.
\end{takeawaybox}
\end{frame}

\begin{frame}{Computer-Use Agents: Strong Benchmarks, Volatile Products}
\textbf{Browser/computer-use agents got good on paper but the standalone products churned.}

\vspace{0.16cm}

\begin{itemize}\footnotesize
    \item Computer-use / browser agents reached $\sim$\textbf{78--84\% on OSWorld-Verified} (human baseline $\sim$72\%).
    \item \textbf{But standalone products are volatile:} Google \emph{killed} Project Mariner (May 2026); OpenAI folded Operator into ChatGPT Agent / Atlas.
\end{itemize}

\vspace{0.12cm}
\begin{cautionbox}
\footnotesize\textbf{Flag the benchmark numbers.} OSWorld scores are \textbf{scaffold-dependent and largely self-reported}: the same model scores very differently under different harnesses, and ``above human baseline'' does not mean reliable on \emph{your} unseen workflow. Treat as directional. \textbf{[solid claim that products churned; benchmark figures = vendor/scaffold-dependent]}
\end{cautionbox}

\vspace{0.12cm}
\textbf{Lesson:} do not build durable research infrastructure on a single experimental product. Build on the protocol layer (MCP) and your own files; treat any one browser agent as replaceable.
\end{frame}

%=============================================================================
\section{Benchmarks \& Measurement Validity}
%===========================================================

\sectiondivider{Benchmarks \& Measurement Validity}{The course's measurement theme, applied to agents}

\begin{frame}{SWE-bench Verified Collapsed Under Scrutiny}
\textbf{This module is the measurement theme of the course pointed at agent evals: a headline number is not a measurement until you know its validity.}

\vspace{0.16cm}

\begin{itemize}\footnotesize
    \item \textbf{Contamination + harness-hacking critiques} eroded SWE-bench Verified; \textbf{OpenAI stopped reporting it in Feb 2026}.
    \item Berkeley RDI \emph{trivially broke four agent benchmarks via the harness} (Apr 2026) --- exploiting the scaffold, not solving the task.
    \item \textbf{Always pair} a Verified headline with \textbf{Scale's SWE-bench Pro} ($\sim$55--59\% frontier, held-out problems).
    \item \textbf{Reliability under repetition} matters more than best-of-$n$ for autonomous use: $\tau$-bench \emph{pass-hat-$k$} falls from $\sim$90\% single-try to $\sim$57\% at $k{=}8$.
\end{itemize}

\vspace{0.1cm}
\begin{cautionbox}
\footnotesize\textbf{Hard rule.} Never put a SWE-bench \emph{Verified} score and a SWE-bench \emph{Pro} score on the same axis --- different problem sets, different difficulty, different contamination exposure. \textbf{[solid]}
\end{cautionbox}
\end{frame}

\begin{frame}{Single-Try Score vs.\ Reliability Under Repetition}
\textbf{For an agent that runs unattended, the relevant statistic is not its \emph{best} run --- it is whether it succeeds \emph{every} time.}

\vspace{0.2cm}

\begin{columns}[T]
\begin{column}{0.5\textwidth}
\textbf{Best-of-$n$ thinking (wrong for autonomy)}
\begin{itemize}\footnotesize
    \item report the best of several attempts
    \item flatters the model
    \item fine for a human-in-the-loop draft you will check
\end{itemize}
\end{column}
\begin{column}{0.48\textwidth}
\textbf{pass-hat-$k$ / reliability (right for autonomy)}
\begin{itemize}\footnotesize
    \item probability \emph{all} $k$ independent runs succeed
    \item $\tau$-bench: $\sim$90\% at $k{=}1$ $\to$ $\sim$57\% at $k{=}8$
    \item this is what matters when no human checks each run
\end{itemize}
\end{column}
\end{columns}

\vspace{0.22cm}

\textbf{Economist translation:} an unattended agent on a recurring data pull is a repeated trial. A 90\% single-try success compounds into frequent failure across a quarter of nightly runs. Design for the reliability curve, not the demo.
\end{frame}

%=============================================================================
\section{Safety \& Reproducibility}
%===========================================================

\sectiondivider{Safety \& Reproducibility}{Untrusted inputs and chain-of-custody for agentic research}

\begin{frame}{Treat Everything the Agent Reads as Untrusted}
\textbf{The defining agentic-AI security failure is prompt injection: data the agent merely \emph{reads} can hijack what it \emph{does}.}

\vspace{0.14cm}

\begin{itemize}\footnotesize
    \item \textbf{OWASP Top-10 for Agentic Applications} (Dec 2025) now codifies these risks.
    \item Documented exploits: \textbf{EchoLeak} (CVE-2025-32711, zero-click M365 Copilot data exfiltration); \textbf{ShadowPrompt} in Claude's Chrome extension; Anthropic's \textbf{GTG-1002} --- the first largely-autonomous cyberattack (Nov 2025).
    \item \textbf{NIST CAISI} red-teaming raised task-hijacking success from \textbf{11\% to 81\%} with adaptive attacks.
\end{itemize}

\vspace{0.12cm}
\begin{cautionbox}
\footnotesize\textbf{Operating rule.} Treat any content the agent reads --- emails, web pages, retrieved documents, even tool descriptions --- as an \emph{untrusted injection channel}. Use least-privilege sandboxing and revocable permissions, exactly as T4's \texttt{settings.json} \texttt{deny} list and T5's security frame prescribed. \textbf{[solid]}
\end{cautionbox}
\end{frame}

\begin{frame}{Reproducibility: Chain-of-Custody for Agentic Results}
\textbf{Can an agent reproduce a published computational result? Often only partly --- which makes logging non-negotiable.}

\vspace{0.16cm}

\begin{itemize}\footnotesize
    \item \textbf{CORE-Bench}: agents reproduce published computational results only $\sim$\textbf{60\% (easy)} / $\sim$\textbf{21\% (hard)}. \textbf{[solid]}
    \item Therefore \textbf{log full trajectories}: prompts, tool calls, seeds, \emph{model versions}, \emph{spec versions}, and \emph{dates} --- the chain-of-custody for any agent-produced artifact.
\end{itemize}

\vspace{0.12cm}
\begin{cautionbox}
\footnotesize\textbf{Read vendor safety numbers skeptically.} Distinguish a vendor's ``$\sim$1\% attack success on a \emph{fixed} suite'' from academic findings that \emph{adaptive} attacks still exceed $\sim$85\%. A fixed-suite number measures the suite, not the threat model. \textbf{[solid]}
\end{cautionbox}

\vspace{0.1cm}
\footnotesize\textbf{Ties to T6.} This is the same six-criteria discipline (reproducibility, verifiability, auditability) --- now with the dates and spec versions the landscape forces you to pin.
\end{frame}

%=============================================================================
\section{Synthesis}
%===========================================================

\begin{frame}{Synthesis: The Landscape Reinforces the Workflow Spine}
\textbf{Every module above maps back onto the terminal-agent workflow of T1--T6 and onto the course's measurement / identification theme.}

\vspace{0.16cm}

\begin{center}
\footnotesize
\renewcommand{\arraystretch}{1.18}
\begin{tabular}{|p{3.2cm}|p{7.2cm}|}
\hline
\textbf{Landscape module} & \textbf{Connection to the spine} \\ \hline
Interoperability (MCP/A2A) & build on the neutral protocol layer, not one product; pin spec versions \\ \hline
Orchestration & T4's sub-agent review \emph{is} orchestrator--worker; pay for multi-agent only when value $>$ token cost \\ \hline
Reasoning \& time horizons & ``runs for hours'' $\to$ a measured, reliability-thresholded number \\ \hline
Coding / computer-use & autonomy is a sandboxing + approval-gate decision, set before the run \\ \hline
Benchmarks & a headline score is not a measurement until you know its validity (Verified vs.\ Pro) \\ \hline
Safety \& reproducibility & untrusted-input discipline $+$ logged chain-of-custody (dates, versions, seeds) \\ \hline
\end{tabular}
\end{center}

\vspace{0.1cm}
\begin{center}
\textbf{The plumbing changes monthly; the research discipline does not.}\\
\footnotesize\textit{Compete on questions and judgment; cooperate on standards; measure before you trust.}
\end{center}
\end{frame}

%===========================================================
\end{document}
